Este capítulo abordará los servicios web necesarios para este proyecto, el modelo de datos, diagramas de entidad-relación y diseño de la plataforma DOMU. Se presenta el modelo de datos, la arquitectura del software y el diseño de las interfaces, que permiten visualizar cómo se implementarán los requerimientos funcionales y no funcionales definidos en el capítulo anterior.

% ---------------------------------------------
\section{Descripción de los Servicios Web}

La capa de servicios web de DOMU se implementa mediante una \textbf{API RESTful} expuesta a través del \textit{framework} Javalin~6, complementada con canales \textbf{WebSocket} para comunicación bidireccional en tiempo real. El backend agrupa la lógica en trece módulos funcionales, cada uno respaldado por sus respectivos \textit{Controllers}, \textit{Services} y \textit{Repositories}. Adicionalmente, el sistema integra servicios transversales de generación de documentos PDF (OpenPDF), almacenamiento en la nube (Box~API) y envío de correo electrónico (Jakarta~Mail).

La Tabla~\ref{tab:servicios-web} resume los módulos, el protocolo de comunicación que utilizan y una descripción sintética de su alcance funcional.

\begin{longtable}{|p{0.22\textwidth}|p{0.13\textwidth}|p{0.57\textwidth}|}
\caption{Módulos de servicios web del sistema DOMU}\label{tab:servicios-web}\\
\hline
\textbf{Módulo} & \textbf{Protocolo} & \textbf{Descripción} \\
\hline
\endfirsthead
\hline
\textbf{Módulo} & \textbf{Protocolo} & \textbf{Descripción} \\
\hline
\endhead
\hline
\endfoot

Autenticación y Autorización &
REST &
Registro con confirmación por correo, inicio de sesión con JWT, recuperación de contraseña mediante token temporal y cierre de sesión. \\
\hline

Gestión de Usuarios y Roles &
REST &
CRUD de usuarios, asignación dinámica de cinco roles (Administrador, Residente, Conserje, Funcionario, Proveedor) bajo el modelo RBAC; gestión de perfiles con avatar almacenado en Box. \\
\hline

Gestión de Comunidades &
REST / E-mail &
Registro de comunidades con aprobación por correo electrónico, invitación de administradores, soporte multi-comunidad y vinculación de usuarios a edificios. \\
\hline

Control de Accesos &
REST &
Registro de visitas, pre-autorización con código QR, check-in/check-out, gestión de contactos frecuentes y bitácora de ingresos. \\
\hline

Gestión de Encomiendas &
REST &
Recepción de paquetes, seguimiento de estados (recibido, notificado, retirado), evidencia fotográfica y notificación al residente. \\
\hline

Gestión Operativa &
REST &
Administración de personal y turnos, creación y asignación de tareas con evidencia fotográfica y seguimiento por estado. \\
\hline

Administración Financiera &
REST / PDF &
Períodos de gasto común, cargos por unidad, registro de pagos, generación de boletas y estados de cuenta en PDF (OpenPDF). \\
\hline

Incidentes y Tickets &
REST &
Reporte de incidencias por residentes, asignación a responsables, seguimiento por estado y prioridad. \\
\hline

Reservas de Áreas Comunes &
REST &
Definición de amenidades con franjas horarias, control de aforo, consulta de disponibilidad y gestión de reservas. \\
\hline

Votaciones &
REST &
Creación de procesos de votación, sufragio remoto, cierre automático y exportación de resultados. \\
\hline

Comunicación Comunitaria &
REST / WebSocket &
Foro con categorías e hilos, chat en tiempo real vía WebSocket (\texttt{/ws/chat}), marketplace de artículos entre vecinos y biblioteca de documentos comunitarios. \\
\hline

Notificaciones &
WebSocket &
Entrega de notificaciones en tiempo real mediante el canal \texttt{/ws/notifications}, con preferencias configurables por tipo de evento. \\
\hline

Gestión de Proveedores y Órdenes de Servicio &
REST &
CRUD de proveedores vinculados a la comunidad, creación y seguimiento de órdenes de servicio con flujo de aceptación/rechazo por parte del proveedor, y gestión de cotizaciones. \\

\end{longtable}

% ---------------------------------------------
\subsection{Especificación de Endpoints Principales}

A continuación se documentan los endpoints REST más representativos del sistema, agrupados por módulo funcional. La API completa cuenta con más de 130~endpoints; para la especificación exhaustiva e interactiva, se recomienda consultar Swagger~UI en la ruta \texttt{/swagger} del servidor.

% --- Endpoints: Autenticación ---
\begin{longtable}{|p{0.08\textwidth}|p{0.38\textwidth}|p{0.46\textwidth}|}
\caption{Endpoints --- Autenticación}\label{tab:ep-auth}\\
\hline
\textbf{Método} & \textbf{Ruta} & \textbf{Descripción} \\
\hline
\endfirsthead
\hline \textbf{Método} & \textbf{Ruta} & \textbf{Descripción} \\ \hline
\endhead
\hline \endfoot

POST & \texttt{/api/auth/register} & Registro de usuario con envío de correo de confirmación. \\
\hline
POST & \texttt{/api/auth/login} & Inicio de sesión; retorna token JWT. \\
\hline
GET & \texttt{/api/auth/confirm-email} & Confirma correo electrónico mediante token. \\
\hline
POST & \texttt{/api/auth/forgot-password} & Envía correo con enlace de recuperación de contraseña. \\
\hline
POST & \texttt{/api/auth/reset-password} & Restablece contraseña mediante token temporal. \\

\end{longtable}

% --- Endpoints: Usuarios ---
\begin{longtable}{|p{0.08\textwidth}|p{0.38\textwidth}|p{0.46\textwidth}|}
\caption{Endpoints --- Usuarios y Perfiles}\label{tab:ep-users}\\
\hline
\textbf{Método} & \textbf{Ruta} & \textbf{Descripción} \\
\hline
\endfirsthead
\hline \textbf{Método} & \textbf{Ruta} & \textbf{Descripción} \\ \hline
\endhead
\hline \endfoot

GET & \texttt{/api/users/profile} & Obtiene perfil del usuario autenticado. \\
\hline
PUT & \texttt{/api/users/profile} & Actualiza datos del perfil. \\
\hline
POST & \texttt{/api/users/avatar} & Sube o actualiza avatar (almacenado en Box). \\
\hline
GET & \texttt{/api/users} & Lista usuarios de la comunidad activa (filtro por rol). \\

\end{longtable}

% --- Endpoints: Comunidades ---
\begin{longtable}{|p{0.08\textwidth}|p{0.38\textwidth}|p{0.46\textwidth}|}
\caption{Endpoints --- Comunidades}\label{tab:ep-communities}\\
\hline
\textbf{Método} & \textbf{Ruta} & \textbf{Descripción} \\
\hline
\endfirsthead
\hline \textbf{Método} & \textbf{Ruta} & \textbf{Descripción} \\ \hline
\endhead
\hline \endfoot

POST & \texttt{/api/buildings} & Solicita registro de nueva comunidad. \\
\hline
POST & \texttt{/api/buildings/approve/\{token\}} & Aprueba solicitud de comunidad vía correo. \\
\hline
POST & \texttt{/api/buildings/reject/\{token\}} & Rechaza solicitud de comunidad. \\
\hline
GET & \texttt{/api/buildings} & Lista comunidades del usuario autenticado. \\
\hline
GET & \texttt{/api/buildings/\{id\}} & Obtiene detalle de una comunidad. \\

\end{longtable}

% --- Endpoints: Control de Accesos ---
\begin{longtable}{|p{0.08\textwidth}|p{0.38\textwidth}|p{0.46\textwidth}|}
\caption{Endpoints --- Control de Accesos}\label{tab:ep-access}\\
\hline
\textbf{Método} & \textbf{Ruta} & \textbf{Descripción} \\
\hline
\endfirsthead
\hline \textbf{Método} & \textbf{Ruta} & \textbf{Descripción} \\ \hline
\endhead
\hline \endfoot

POST & \texttt{/api/visits} & Registra nueva visita (ingreso). \\
\hline
POST & \texttt{/api/visits/qr-check} & Verifica código QR de visita pre-autorizada. \\
\hline
PUT & \texttt{/api/visits/\{id\}/checkout} & Registra salida de visita. \\
\hline
GET & \texttt{/api/visits} & Lista visitas de la comunidad (filtros por fecha y estado). \\
\hline
POST & \texttt{/api/visits/pre-authorize} & Pre-autoriza visita con generación de QR. \\
\hline
GET & \texttt{/api/visit-contacts} & Lista contactos frecuentes del residente. \\

\end{longtable}

% --- Endpoints: Encomiendas ---
\begin{longtable}{|p{0.08\textwidth}|p{0.38\textwidth}|p{0.46\textwidth}|}
\caption{Endpoints --- Encomiendas}\label{tab:ep-parcels}\\
\hline
\textbf{Método} & \textbf{Ruta} & \textbf{Descripción} \\
\hline
\endfirsthead
\hline \textbf{Método} & \textbf{Ruta} & \textbf{Descripción} \\ \hline
\endhead
\hline \endfoot

POST & \texttt{/api/parcels} & Registra nueva encomienda con evidencia fotográfica. \\
\hline
PUT & \texttt{/api/parcels/\{id\}/status} & Actualiza estado (notificado, retirado). \\
\hline
GET & \texttt{/api/parcels} & Lista encomiendas de la comunidad. \\
\hline
GET & \texttt{/api/parcels/\{id\}} & Detalle de una encomienda. \\

\end{longtable}

% --- Endpoints: Finanzas ---
\begin{longtable}{|p{0.08\textwidth}|p{0.38\textwidth}|p{0.46\textwidth}|}
\caption{Endpoints --- Administración Financiera}\label{tab:ep-finance}\\
\hline
\textbf{Método} & \textbf{Ruta} & \textbf{Descripción} \\
\hline
\endfirsthead
\hline \textbf{Método} & \textbf{Ruta} & \textbf{Descripción} \\ \hline
\endhead
\hline \endfoot

POST & \texttt{/api/common-expenses/periods} & Crea nuevo período de gasto común. \\
\hline
POST & \texttt{/api/common-expenses/charges} & Agrega cargos a un período por unidad. \\
\hline
POST & \texttt{/api/payments} & Registra pago de gasto común. \\
\hline
GET & \texttt{/api/common-expenses/statement/\{unitId\}} & Obtiene estado de cuenta de una unidad. \\
\hline
GET & \texttt{/api/common-expenses/statement/\{unitId\}/pdf} & Descarga estado de cuenta en PDF. \\
\hline
GET & \texttt{/api/payments} & Lista pagos registrados (filtros por período y unidad). \\

\end{longtable}

% --- Endpoints: Incidentes ---
\begin{longtable}{|p{0.08\textwidth}|p{0.38\textwidth}|p{0.46\textwidth}|}
\caption{Endpoints --- Incidentes y Tickets}\label{tab:ep-incidents}\\
\hline
\textbf{Método} & \textbf{Ruta} & \textbf{Descripción} \\
\hline
\endfirsthead
\hline \textbf{Método} & \textbf{Ruta} & \textbf{Descripción} \\ \hline
\endhead
\hline \endfoot

POST & \texttt{/api/incidents} & Crea nuevo ticket de incidencia. \\
\hline
PUT & \texttt{/api/incidents/\{id\}/status} & Actualiza estado del ticket. \\
\hline
GET & \texttt{/api/incidents} & Lista incidentes de la comunidad. \\
\hline
GET & \texttt{/api/incidents/stats} & Obtiene estadísticas de incidentes. \\

\end{longtable}

% --- Endpoints: Reservas ---
\begin{longtable}{|p{0.08\textwidth}|p{0.38\textwidth}|p{0.46\textwidth}|}
\caption{Endpoints --- Reservas de Áreas Comunes}\label{tab:ep-reservations}\\
\hline
\textbf{Método} & \textbf{Ruta} & \textbf{Descripción} \\
\hline
\endfirsthead
\hline \textbf{Método} & \textbf{Ruta} & \textbf{Descripción} \\ \hline
\endhead
\hline \endfoot

GET & \texttt{/api/amenities} & Lista áreas comunes de la comunidad. \\
\hline
GET & \texttt{/api/amenities/\{id\}/availability} & Consulta disponibilidad por fecha. \\
\hline
POST & \texttt{/api/reservations} & Crea reserva de área común. \\
\hline
DELETE & \texttt{/api/reservations/\{id\}} & Cancela reserva existente. \\

\end{longtable}

% --- Endpoints: Votaciones ---
\begin{longtable}{|p{0.08\textwidth}|p{0.38\textwidth}|p{0.46\textwidth}|}
\caption{Endpoints --- Votaciones}\label{tab:ep-polls}\\
\hline
\textbf{Método} & \textbf{Ruta} & \textbf{Descripción} \\
\hline
\endfirsthead
\hline \textbf{Método} & \textbf{Ruta} & \textbf{Descripción} \\ \hline
\endhead
\hline \endfoot

POST & \texttt{/api/polls} & Crea nueva votación con opciones y fecha de cierre. \\
\hline
POST & \texttt{/api/polls/\{id\}/vote} & Registra voto del usuario autenticado. \\
\hline
PUT & \texttt{/api/polls/\{id\}/close} & Cierra votación y consolida resultados. \\
\hline
GET & \texttt{/api/polls/\{id\}/results} & Obtiene resultados de la votación. \\

\end{longtable}

% --- Endpoints: Chat ---
\begin{longtable}{|p{0.08\textwidth}|p{0.38\textwidth}|p{0.46\textwidth}|}
\caption{Endpoints --- Chat}\label{tab:ep-chat}\\
\hline
\textbf{Método} & \textbf{Ruta} & \textbf{Descripción} \\
\hline
\endfirsthead
\hline \textbf{Método} & \textbf{Ruta} & \textbf{Descripción} \\ \hline
\endhead
\hline \endfoot

POST & \texttt{/api/chat/rooms} & Crea sala de chat o solicita conversación. \\
\hline
POST & \texttt{/api/chat/rooms/\{id\}/messages} & Envía mensaje a una sala (REST fallback). \\
\hline
GET & \texttt{/api/chat/rooms/\{id\}/messages} & Obtiene historial de mensajes. \\

\end{longtable}

% --- Endpoints: Notificaciones ---
\begin{longtable}{|p{0.08\textwidth}|p{0.38\textwidth}|p{0.46\textwidth}|}
\caption{Endpoints --- Notificaciones}\label{tab:ep-notifications}\\
\hline
\textbf{Método} & \textbf{Ruta} & \textbf{Descripción} \\
\hline
\endfirsthead
\hline \textbf{Método} & \textbf{Ruta} & \textbf{Descripción} \\ \hline
\endhead
\hline \endfoot

GET & \texttt{/api/notifications} & Lista notificaciones del usuario. \\
\hline
PUT & \texttt{/api/notifications/\{id\}/read} & Marca notificación como leída. \\
\hline
GET & \texttt{/api/notifications/preferences} & Obtiene preferencias de notificación. \\

\end{longtable}

% --- Endpoints: Proveedores ---
\begin{longtable}{|p{0.08\textwidth}|p{0.38\textwidth}|p{0.46\textwidth}|}
\caption{Endpoints --- Proveedores}\label{tab:ep-providers}\\
\hline
\textbf{Método} & \textbf{Ruta} & \textbf{Descripción} \\
\hline
\endfirsthead
\hline \textbf{Método} & \textbf{Ruta} & \textbf{Descripción} \\ \hline
\endhead
\hline \endfoot

POST & \texttt{/api/providers} & Registra nuevo proveedor en la comunidad. \\
\hline
GET & \texttt{/api/providers} & Lista proveedores de la comunidad. \\
\hline
GET & \texttt{/api/providers/\{id\}} & Detalle de un proveedor. \\
\hline
PUT & \texttt{/api/providers/\{id\}} & Actualiza datos del proveedor. \\
\hline
DELETE & \texttt{/api/providers/\{id\}} & Desactiva proveedor. \\

\end{longtable}

% --- Endpoints: Órdenes de Servicio ---
\begin{longtable}{|p{0.08\textwidth}|p{0.38\textwidth}|p{0.46\textwidth}|}
\caption{Endpoints --- Órdenes de Servicio}\label{tab:ep-service-orders}\\
\hline
\textbf{Método} & \textbf{Ruta} & \textbf{Descripción} \\
\hline
\endfirsthead
\hline \textbf{Método} & \textbf{Ruta} & \textbf{Descripción} \\ \hline
\endhead
\hline \endfoot

POST & \texttt{/api/service-orders} & Crea orden de servicio asignada a proveedor. \\
\hline
PUT & \texttt{/api/service-orders/\{id\}/status} & Actualiza estado de la orden (aceptar, rechazar, completar). \\
\hline
POST & \texttt{/api/service-orders/\{id\}/quotations} & Sube cotización vinculada a la orden. \\
\hline
GET & \texttt{/api/service-orders} & Lista órdenes de servicio (filtros por estado y proveedor). \\

\end{longtable}

Adicionalmente, el sistema expone dos canales \textbf{WebSocket} persistentes:
\begin{itemize}
    \item \texttt{/ws/chat} --- Mensajería en tiempo real entre usuarios de la comunidad.
    \item \texttt{/ws/notifications} --- Entrega de notificaciones en tiempo real según preferencias del usuario.
\end{itemize}

Ambos canales requieren autenticación mediante token JWT enviado como parámetro en la URL de conexión.

% === PLACEHOLDER: Reemplazar con imagen generada ===
% PROMPT para generador de imágenes (Gemini/GPT/DALL-E):
%
% "Genera un DIAGRAMA DE COMPONENTES/MÓDULOS de una API REST para un sistema llamado DOMU.
%
% NOTACIÓN:
% - MÓDULOS: rectángulos redondeados con borde sólido, nombre del módulo dentro en negrita.
% - CORE CENTRAL: un rectángulo grande en el centro etiquetado 'API REST - Javalin 6' con ícono de servidor.
% - WEBSOCKET: 2 rectángulos con borde punteado y doble flecha bidireccional, etiquetados '/ws/chat' y '/ws/notifications'.
% - CLIENTE: a la izquierda, un rectángulo etiquetado 'Frontend React 19 (SPA)' con flecha HTTP hacia el core.
% - BASE DE DATOS: a la derecha, un cilindro etiquetado 'MySQL 8' con flecha desde el core.
% - SERVICIOS EXTERNOS: arriba, 3 cajas pequeñas: 'Box API (archivos)', 'Gmail SMTP (correo)', 'OpenPDF (documentos)'.
%
% Los 13 módulos se organizan en GRID 4x4 alrededor del core central:
% Fila superior: Autenticación, Usuarios y Roles, Comunidades, Control de Accesos
% Fila media: Encomiendas, Gestión Operativa, Finanzas, Incidentes
% Fila inferior: Reservas, Votaciones, Comunicación (Foro/Chat/Marketplace), Notificaciones
% Última fila: Proveedores y Órdenes de Servicio
%
% Cada módulo tiene una línea que lo conecta al core central.
% Los módulos de Comunicación y Notificaciones tienen además línea punteada a los WebSocket.
%
% ESTILO: Fondo blanco, bordes negros o gris oscuro, relleno de módulos en gris muy claro (#F5F5F5).
% Sin gradientes, sin sombras. Tipografía sans-serif. Profesional para tesis universitaria.
% Etiquetas en español. Formato: PNG 300dpi."
\begin{figure}[H]
  \centering
  \fbox{\parbox{0.85\textwidth}{\centering\vspace{3cm}
  \textit{[Insertar diagrama: Arquitectura modular de la API REST --- 13 módulos + 2 canales WebSocket]}
  \vspace{3cm}}}
  \caption{Diagrama de módulos de la API REST del sistema DOMU.}
  \label{fig:diagrama-api-modulos}
\end{figure}

% ---------------------------------------------
\section{Modelo de datos}

El modelo de datos es de tipo relacional, con el fin de mantener la simpleza y eficiencia del almacenamiento de información crítica del sistema.

Las 37~tablas del modelo, derivadas de las 35~migraciones Flyway del proyecto, se organizan en 12~dominios funcionales. La Tabla~\ref{tab:resumen-dominios} presenta un resumen de esta distribución; el detalle de cada entidad y sus atributos se encuentra en el Anexo~D (Diccionario de Datos).

\begin{table}[H]
\centering
\small
\begin{tabular}{|p{0.28\textwidth}|p{0.52\textwidth}|c|}
\hline
\textbf{Dominio} & \textbf{Entidades principales} & \textbf{Cant.} \\ \hline
Identidad y Roles & \texttt{users}, \texttt{roles}, \texttt{user\_confirmations}, \texttt{password\_reset\_tokens} & 4 \\ \hline
Comunidad y Estructura & \texttt{buildings}, \texttt{housing\_units}, \texttt{user\_buildings} & 3 \\ \hline
Solicitudes de Comunidad & \texttt{building\_requests} & 1 \\ \hline
Visitas y Control de Acceso & \texttt{visits}, \texttt{visit\_authorizations}, \texttt{visit\_access\_logs}, \texttt{visit\_contacts} & 4 \\ \hline
Gestión Operativa & \texttt{staff}, \texttt{tasks}, \texttt{task\_assignments}, \texttt{parcels} & 4 \\ \hline
Finanzas y Gastos Comunes & \texttt{common\_expense\_periods}, \texttt{common\_charges}, \texttt{common\_payments}, \texttt{delinquency\_records}, \texttt{common\_expense\_revisions} & 5 \\ \hline
Reservas de Espacios & \texttt{amenities}, \texttt{amenity\_time\_slots}, \texttt{amenity\_reservations} & 3 \\ \hline
Votaciones & \texttt{polls}, \texttt{poll\_options}, \texttt{poll\_votes} & 3 \\ \hline
Chat y Marketplace & \texttt{chat\_room}, \texttt{chat\_participant}, \texttt{chat\_message}, \texttt{chat\_request}, \texttt{market\_category}, \texttt{market\_item}, \texttt{market\_item\_image} & 7 \\ \hline
Foro Comunitario & \texttt{forum\_categories}, \texttt{forum\_threads}, \texttt{forum\_posts} & 3 \\ \hline
Biblioteca, Notificaciones e Incidentes & \texttt{library\_documents}, \texttt{notifications}, \texttt{notification\_preferences}, \texttt{incidents} & 4 \\ \hline
Proveedores & \texttt{providers}, \texttt{service\_orders}, \texttt{quotations} & 3 \\ \hline
\end{tabular}
\caption{Resumen de entidades del modelo de datos por dominio funcional.}
\label{tab:resumen-dominios}
\end{table}

\subsection{Esquema de la base de datos}
El diseño de la persistencia de datos se realizó utilizando el motor de base de datos relacional \textbf{MySQL~8}. El esquema se gestiona mediante \textbf{Flyway~10.15}, que ejecuta migraciones SQL versionadas al arranque del servidor, garantizando la consistencia entre los entornos de desarrollo y producción. Hasta la fecha, el proyecto acumula \textbf{35~migraciones} (\texttt{V001} a \texttt{V035}), cada una atómica e idempotente.

A continuación, las Tablas~\ref{tab:esquema-acceso} a~\ref{tab:esquema-proveedores} presentan las tablas del esquema agrupadas por dominio funcional, junto con sus relaciones clave.

% --- Tabla 1: Acceso y Organización ---
\begin{table}[H]
\centering
\caption{Esquema de base de datos --- Acceso y Organización}\label{tab:esquema-acceso}
\begin{tabular}{|p{0.26\textwidth}|p{0.38\textwidth}|p{0.28\textwidth}|}
\hline
\textbf{Tabla} & \textbf{Descripción} & \textbf{Relaciones clave} \\
\hline
\texttt{roles} & Catálogo de cinco roles del sistema. & --- \\
\hline
\texttt{users} & Credenciales, datos personales y avatar. & FK a \texttt{roles} \\
\hline
\texttt{buildings} & Comunidades con dirección y configuración. & --- \\
\hline
\texttt{housing\_units} & Unidades habitacionales con prorrata. & FK a \texttt{buildings} \\
\hline
\texttt{user\_buildings} & Vínculo usuarios--comunidades (N:M). & FK a \texttt{users}, \texttt{buildings} \\
\hline
\texttt{user\_confirmations} & Tokens de verificación de correo. & FK a \texttt{users} \\
\hline
\texttt{password\_reset\_tokens} & Tokens temporales de recuperación. & FK a \texttt{users} \\
\hline
\end{tabular}
\end{table}

% --- Tabla 2: Solicitudes de Comunidad y Control de Acceso ---
\begin{table}[H]
\centering
\caption{Esquema de base de datos --- Solicitudes de Comunidad y Control de Acceso}\label{tab:esquema-solicitudes}
\begin{tabular}{|p{0.26\textwidth}|p{0.38\textwidth}|p{0.28\textwidth}|}
\hline
\textbf{Tabla} & \textbf{Descripción} & \textbf{Relaciones clave} \\
\hline
\texttt{building\_requests} & Solicitudes de ingreso con aprobación. & FK a \texttt{buildings}, \texttt{users} \\
\hline
\texttt{admin\_invites} & Códigos de invitación para administradores. & FK a \texttt{buildings} \\
\hline
\texttt{visits} & Registros de ingreso de visitantes. & FK a \texttt{buildings}, \texttt{users} \\
\hline
\texttt{visit\_contacts} & Contactos frecuentes pre-autorizados. & FK a \texttt{users} \\
\hline
\end{tabular}
\end{table}

% --- Tabla 3: Gestión Operativa ---
\begin{table}[H]
\centering
\caption{Esquema de base de datos --- Gestión Operativa}\label{tab:esquema-operativa}
\begin{tabular}{|p{0.26\textwidth}|p{0.38\textwidth}|p{0.28\textwidth}|}
\hline
\textbf{Tabla} & \textbf{Descripción} & \textbf{Relaciones clave} \\
\hline
\texttt{staff} & Personal de la comunidad. & FK a \texttt{buildings} \\
\hline
\texttt{tasks} & Tareas con evidencia y estados. & FK a \texttt{buildings} \\
\hline
\texttt{task\_assignments} & Asignación tareas--personal. & FK a \texttt{tasks}, \texttt{staff} \\
\hline
\texttt{parcels} & Encomiendas con seguimiento de estados. & FK a \texttt{buildings}, \texttt{users} \\
\hline
\end{tabular}
\end{table}

% --- Tabla 4: Finanzas ---
\begin{table}[H]
\centering
\caption{Esquema de base de datos --- Finanzas}\label{tab:esquema-finanzas}
\begin{tabular}{|p{0.26\textwidth}|p{0.38\textwidth}|p{0.28\textwidth}|}
\hline
\textbf{Tabla} & \textbf{Descripción} & \textbf{Relaciones clave} \\
\hline
\texttt{common\_expense\_periods} & Períodos contables mensuales. & FK a \texttt{buildings} \\
\hline
\texttt{charges} & Cargos individuales por unidad. & FK a \texttt{common\_expense\_periods} \\
\hline
\texttt{payments} & Pagos registrados por unidad. & FK a \texttt{charges}, \texttt{users} \\
\hline
\end{tabular}
\end{table}

% --- Tabla 5: Participación ---
\begin{table}[H]
\centering
\caption{Esquema de base de datos --- Participación y Reservas}\label{tab:esquema-participacion}
\begin{tabular}{|p{0.26\textwidth}|p{0.38\textwidth}|p{0.28\textwidth}|}
\hline
\textbf{Tabla} & \textbf{Descripción} & \textbf{Relaciones clave} \\
\hline
\texttt{polls} & Procesos de votación comunitaria. & FK a \texttt{buildings} \\
\hline
\texttt{poll\_options} & Alternativas de cada votación. & FK a \texttt{polls} \\
\hline
\texttt{votes} & Sufragios individuales. & FK a \texttt{poll\_options}, \texttt{users} \\
\hline
\texttt{amenities} & Espacios comunes con reglas de aforo. & FK a \texttt{buildings} \\
\hline
\texttt{reservations} & Reservas de bloques horarios. & FK a \texttt{amenities}, \texttt{users} \\
\hline
\end{tabular}
\end{table}

% --- Tabla 6: Comunicación ---
\begin{table}[H]
\centering
\caption{Esquema de base de datos --- Comunicación}\label{tab:esquema-comunicacion}
\begin{tabular}{|p{0.26\textwidth}|p{0.38\textwidth}|p{0.28\textwidth}|}
\hline
\textbf{Tabla} & \textbf{Descripción} & \textbf{Relaciones clave} \\
\hline
\texttt{chat\_rooms} & Salas de conversación entre usuarios. & FK a \texttt{users} \\
\hline
\texttt{chat\_messages} & Mensajes individuales con marca temporal. & FK a \texttt{chat\_rooms}, \texttt{users} \\
\hline
\texttt{chat\_requests} & Solicitudes de inicio de conversación. & FK a \texttt{users} \\
\hline
\texttt{marketplace\_items} & Artículos en venta o donación. & FK a \texttt{users}, \texttt{buildings} \\
\hline
\texttt{forum\_categories} & Categorías temáticas del foro. & FK a \texttt{buildings} \\
\hline
\texttt{forum\_threads} & Hilos de discusión. & FK a \texttt{forum\_categories}, \texttt{users} \\
\hline
\end{tabular}
\end{table}

% --- Tabla 7: Soporte ---
\begin{table}[H]
\centering
\caption{Esquema de base de datos --- Soporte}\label{tab:esquema-soporte}
\begin{tabular}{|p{0.26\textwidth}|p{0.38\textwidth}|p{0.28\textwidth}|}
\hline
\textbf{Tabla} & \textbf{Descripción} & \textbf{Relaciones clave} \\
\hline
\texttt{incidents} & Reportes de incidencias. & FK a \texttt{buildings}, \texttt{users} \\
\hline
\texttt{library\_documents} & Documentos comunitarios en Box. & FK a \texttt{buildings} \\
\hline
\texttt{notifications} & Alertas entregadas vía WebSocket. & FK a \texttt{users}, \texttt{buildings} \\
\hline
\texttt{notification\_preferences} & Preferencias por tipo de evento. & FK a \texttt{users} \\
\hline
\end{tabular}
\end{table}

% --- Tabla 8: Proveedores ---
\begin{table}[H]
\centering
\caption{Esquema de base de datos --- Proveedores y Servicios Externos}\label{tab:esquema-proveedores}
\begin{tabular}{|p{0.26\textwidth}|p{0.38\textwidth}|p{0.28\textwidth}|}
\hline
\textbf{Tabla} & \textbf{Descripción} & \textbf{Relaciones clave} \\
\hline
\texttt{providers} & Proveedores externos de servicios. & FK a \texttt{buildings}, \texttt{users} \\
\hline
\texttt{service\_orders} & Órdenes de servicio con flujo de estados. & FK a \texttt{buildings}, \texttt{providers}, \texttt{users} \\
\hline
\texttt{quotations} & Cotizaciones asociadas a órdenes. & FK a \texttt{service\_orders}, \texttt{providers} \\
\hline
\end{tabular}
\end{table}

El modelo respeta las formas normales para evitar redundancia de datos, utilizando claves foráneas (\texttt{FOREIGN KEY}) para mantener la integridad referencial entre las unidades, los residentes y sus respectivas operaciones (pagos, reservas, incidentes). Todas las tablas incluyen columnas de auditoría \texttt{created\_at} y \texttt{updated\_at} gestionadas automáticamente.

\cleardoublepage
\subsection{Modelo Entidad-Relación}

Ya que el proyecto utilizará MySQL como base de datos, definimos el siguiente diagrama Entidad-Relación que ayudará a comprender mejor la definición de las relaciones y la estructura de los datos.

\begin{figure}[H]
  \centering
  \includegraphics[width=\textwidth]{MER.png}
  \caption{Modelo Entidad-Relación del Sistema DOMU}
  \label{fig:mer}
\end{figure}

\clearpage

\cleardoublepage
\subsection{Modelo Relacional}

A continuación, se modelaron los datos siguiendo el patrón relacional, detallando las tablas, columnas, claves primarias y foráneas que se implementarán en MySQL.

\begin{figure}[H]
  \centering
  \includegraphics[width=\textwidth]{ModeloRelacional.png}
  \caption{Diagrama del Modelo Relacional del Sistema DOMU}
  \label{fig:relacional}
\end{figure}

\clearpage

\section{Diseño de interfaz y navegación}
En esta sección se dará énfasis al desarrollo de los \textit{mockups} que contempla nuestra aplicación \textbf{DOMU}, donde se detalla la guía de estilo a seguir, en conjunto con el logotipo que tendrá la \textit{app}. También se definirán las guías de colores a utilizar, las cuales serán las encargadas de representar nuestra identidad visual a través de las interfaces gráficas de la aplicación.

\subsection{Guías de estilos}

La guía de estilo marcará las pautas a seguir para el diseño de la web. Por tanto, servirá de consulta para visualizar los objetivos de la aplicación en cuanto a su estilo. El principal objetivo es dotar al estilo de la aplicación de la máxima sencillez posible, a fin de que al usuario le resulte intuitivo el uso de la aplicación. A continuación se señalarán aspectos relativos al estilo como son el logotipo, los colores principales y secundarios, la tipografía y la composición de las interfaces.

\subsection{Logotipo}
\label{subsec:logotipo}

El logotipo de \textbf{DOMU} es el principal símbolo gráfico que representa la plataforma. Podrá presentarse como \textbf{imagotipo} e \textbf{isotipo}; por defecto se utiliza exclusivamente el imagotipo a color, que combina el ícono con el nombre de la plataforma.

El ícono está basado en una letra “D” tridimensional, diseñada con geometría simple; el contorno evoca el entorno habitacional y comunitario en el que se desarrolla la aplicación.

El texto que acompaña al ícono utiliza una tipografía de estilo moderno y geométrico que refuerza la estética tecnológica de la marca.

\begin{figure}[H]
  \centering
  
\includegraphics[scale=0.5]{LogotipoDOMU.png}
  \caption{Logotipo de Domu}
  \label{fig:logo}
\end{figure}

\subsection{Guía de colores}
Los colores corporativos de la aplicación componen una gama cromática cálida, orientada a transmitir cercanía, energía y modernidad.

En conjunto, esta paleta de colores, que se muestra en la figura a continuación, busca lograr una experiencia visual agradable y llamativa para los usuarios de la plataforma.

\begin{figure}[H]
  \centering
  \includegraphics[scale=0.5]{PaletaDeColores.png}
  \caption{Paleta de Colores DOMU}
  \label{fig:paleta}
\end{figure}

\subsection{Tipografía}
La tipografía utilizada es un tipo de letra sencilla y universal que facilite la legibilidad. Se ha optado por emplear una fuente gratuita de Google Fonts, como es Roboto Mono. En las figuras \ref{fig:tipografia1} y \ref{fig:tipografia2} se incluye una representación de dicha fuente con los caracteres más comunes.

\begin{figure}[H]
  \centering
  \includegraphics[scale=0.5]{EjemploTipografiaRoboto1.png}
  \caption{Texto de Ejemplo 1 Tipografía Roboto Mono}
  \label{fig:tipografia1}
\end{figure}

\begin{figure}[H]
  \centering
  \includegraphics[scale=0.5]{EjemploTipografiaRoboto2.png}
  \caption{Texto de Ejemplo 2 Tipografía Roboto Mono}
  \label{fig:tipografia2}
\end{figure}

\subsection{Composición de Interfaces}
\label{subsec:composicion-interfaces}

La composición de las interfaces expuesta a continuación corresponde a la versión adaptada para móvil que será la versión implementada en primer lugar. La adaptación para pantallas de mayor tamaño se realizará en un futuro.

% --- ÍTEM A: Inicio de Sesión ---
\begin{itemize}
    \item[\textbf{A.}] \textbf{Inicio de Sesión (Login)}: Interfaz de autenticación para todos los usuarios.
\end{itemize}
\begin{figure}[H]
    \centering
    \includegraphics[scale=0.5]{InterfazInicioSesion.png}
    \caption{A. Mockup de la interfaz de Inicio de Sesión.}
    \label{fig:mockup-login}
\end{figure}
\clearpage


% --- ÍTEM B: Vista Conserje (Registro QR) ---
\begin{itemize}
    \item[\textbf{B.}] \textbf{Vista Conserje}: Interfaz operativa clave para el registro de visitas mediante escaneo QR.
\end{itemize}
\begin{figure}[H]
    \centering
    \includegraphics[scale=0.5]{InterfazConserje.png}
    \caption{B. Mockup de la interfaz de la Vista Conserje.}
    \label{fig:mockup-conserje}
\end{figure}
\clearpage


% --- ÍTEM C: Vista Administración ---
\begin{itemize}
    \item[\textbf{C.}] \textbf{Vista Administración}: Paneles de control para la gestión general de la comunidad.
\end{itemize}
\begin{figure}[H]
    \centering
    \includegraphics[scale=0.5]{InterfazAdministracion.png}
    \caption{C. Mockup de la interfaz de la Vista Administración.}
    \label{fig:mockup-administracion}
\end{figure}
\clearpage

% --- ÍTEM D: Vista Comité ---
\begin{itemize}
    \item[\textbf{D.}] \textbf{Vista Comité}: Interfaz con acceso a información crítica, documentos y procesos de votación destinados a los miembros del comité de administración.
\end{itemize}
\begin{figure}[H]
    \centering
    \includegraphics[scale=0.5]{InterfazComite.png}
    \caption{D. Mockup de la interfaz de la Vista Comité.}
    \label{fig:mockup-comite}
\end{figure}
\clearpage


% --- ÍTEM E: Vista Residente (Aplicación Móvil) ---
\begin{itemize}
    \item[\textbf{E.}] \textbf{Vista Residente}: Interfaz de la aplicación móvil, centralizando funcionalidades como recepción de notificaciones, reserva de espacios y pago de gastos comunes.
\end{itemize}
\begin{figure}[H]
    \centering
    \includegraphics[scale=0.5]{InterfazResidente.png}
    \caption{E. Mockup de la interfaz de la Vista Residente (Aplicación Móvil).}
    \label{fig:mockup-residente}
\end{figure}
\clearpage


% --- ÍTEM F: Vista Reporte de Incidencias ---
\begin{itemize}
    \item[\textbf{F.}] \textbf{Vista Reporte de Incidencias}: Interfaz para la gestión y el seguimiento de reclamos, sugerencias o fallas operacionales (\textit{tickets}).
\end{itemize}
\begin{figure}[H]
    \centering
    \includegraphics[scale=0.5]{InterfazReporteDeIncidencias.png}
    \caption{F. Mockup de la interfaz de Reporte de Incidencias.}
    \label{fig:mockup-incidencias}
\end{figure}
\clearpage


% --- ÍTEM G: Vista Reporte Financiero ---
\begin{itemize}
    \item[\textbf{G.}] \textbf{Vista Reporte Financiero}: Paneles analíticos que presentan a la administración las métricas clave de gastos comunes, egresos y estado de morosidad.
\end{itemize}
\begin{figure}[H]
    \centering
    \includegraphics[scale=0.5]{InterfazReportesFinancieros.png}
    \caption{G. Mockup de la interfaz de Reporte Financiero.}
    \label{fig:mockup-financiero}
\end{figure}
\clearpage


% --- ÍTEM H: Vista Pago de Gastos Comunes ---
\begin{itemize}
    \item[\textbf{H.}] \textbf{Vista Pago de Gastos Comunes}: Flujo transaccional que guía al residente a través del proceso de pago en línea de su cuenta mensual.
\end{itemize}
\begin{figure}[H]
    \centering
    \includegraphics[scale=0.5]{InterfazPagoGGCC.png}
    \caption{H. Mockup de la interfaz de Pago de Gastos Comunes.}
    \label{fig:mockup-pago-gc}
\end{figure}
\clearpage

% --- ÍTEM I: Vista Reserva de Áreas Comunes ---
\begin{itemize}
    \item[\textbf{I.}] \textbf{Vista Reserva de Áreas Comunes}: Calendario y gestión de reservas por parte del residente.
\end{itemize}
\begin{figure}[H]
    \centering
    \includegraphics[scale=0.5]{InterfazReservaAreasComunes.png}
    \caption{I. Mockup de la interfaz de Reserva de Áreas Comunes.}
    \label{fig:mockup-reservas}
\end{figure}
\clearpage

\subsection{Vistas del Sistema por Rol}
\label{subsec:vistas-por-rol}

El frontend de DOMU está compuesto por \textbf{54~vistas} organizadas según el rol del usuario. A continuación se presentan las vistas agrupadas por categoría de acceso, indicando los casos de uso (CU) que cada vista materializa.

% --- Vistas Públicas ---
\begin{longtable}{|p{0.28\textwidth}|p{0.42\textwidth}|p{0.22\textwidth}|}
\caption{Vistas públicas (sin autenticación)}\label{tab:vistas-publicas}\\
\hline
\textbf{Vista} & \textbf{Descripción} & \textbf{CUs relacionados} \\
\hline
\endfirsthead
\hline \textbf{Vista} & \textbf{Descripción} & \textbf{CUs relacionados} \\ \hline
\endhead
\hline \endfoot

Home (Landing) & Página de presentación de la plataforma DOMU. & --- \\
\hline
Login & Formulario de inicio de sesión con JWT. & --- \\
\hline
Register & Registro de nuevo usuario con confirmación por correo. & --- \\
\hline
ForgotPassword & Solicitud de recuperación de contraseña. & --- \\
\hline
ConfirmEmail & Confirmación de correo electrónico vía token. & --- \\
\hline
About & Información sobre el proyecto y el equipo. & --- \\
\hline
RegisterByInvite & Registro mediante código de invitación de administrador. & CU\_24 \\

\end{longtable}

% --- Vistas Administrador ---
\begin{longtable}{|p{0.28\textwidth}|p{0.42\textwidth}|p{0.22\textwidth}|}
\caption{Vistas del Administrador}\label{tab:vistas-admin}\\
\hline
\textbf{Vista} & \textbf{Descripción} & \textbf{CUs relacionados} \\
\hline
\endfirsthead
\hline \textbf{Vista} & \textbf{Descripción} & \textbf{CUs relacionados} \\ \hline
\endhead
\hline \endfoot

AdminDashboard & Panel principal con resumen de actividad comunitaria. & CU\_09 \\
\hline
AdminCreateUser & Creación e invitación de usuarios a la comunidad. & CU\_14, CU\_24 \\
\hline
AdminIncidentsBoard & Tablero Kanban de incidentes con estados de avance. & CU\_17 \\
\hline
AdminIncidentStats & Estadísticas y métricas de incidentes reportados. & CU\_09 \\
\hline
AdminResidents & Listado y gestión de residentes por unidad. & CU\_14 \\
\hline
AdminHousingUnits & Gestión de unidades habitacionales y coeficientes. & CU\_24 \\
\hline
AdminAmenities & Configuración de áreas comunes, aforo y horarios. & CU\_05 \\
\hline
AdminCommonExpenses & Gestión de períodos, cargos y pagos de gastos comunes. & CU\_04, CU\_09 \\
\hline
AdminParcels & Gestión y seguimiento de encomiendas. & CU\_02 \\
\hline
AdminTasks & Creación y asignación de tareas al personal. & CU\_07 \\
\hline
AdminStaff & Gestión de personal operativo de la comunidad. & CU\_07 \\
\hline
AdminProviders & CRUD de proveedores vinculados a la comunidad. & CU\_19, CU\_25 \\
\hline
AdminServiceOrders & Gestión de órdenes de servicio y cotizaciones. & CU\_26 \\

\end{longtable}

% --- Vistas Conserje ---
\begin{longtable}{|p{0.28\textwidth}|p{0.42\textwidth}|p{0.22\textwidth}|}
\caption{Vistas del Conserje}\label{tab:vistas-conserje}\\
\hline
\textbf{Vista} & \textbf{Descripción} & \textbf{CUs relacionados} \\
\hline
\endfirsthead
\hline \textbf{Vista} & \textbf{Descripción} & \textbf{CUs relacionados} \\ \hline
\endhead
\hline \endfoot

VisitRegistrationPanel & Registro de visitas con escaneo QR y check-in/out. & CU\_01, CU\_15 \\
\hline
ConciergeVisitLog & Bitácora de visitas registradas. & CU\_01, CU\_03 \\
\hline
ConciergeParcels & Registro y gestión de encomiendas recibidas. & CU\_02 \\

\end{longtable}

% --- Vistas Residente ---
\begin{longtable}{|p{0.28\textwidth}|p{0.42\textwidth}|p{0.22\textwidth}|}
\caption{Vistas del Residente}\label{tab:vistas-residente}\\
\hline
\textbf{Vista} & \textbf{Descripción} & \textbf{CUs relacionados} \\
\hline
\endfirsthead
\hline \textbf{Vista} & \textbf{Descripción} & \textbf{CUs relacionados} \\ \hline
\endhead
\hline \endfoot

ResidentProfile & Perfil personal con edición de datos y avatar. & --- \\
\hline
ResidentVisits & Pre-autorización de visitas con generación de QR. & CU\_15 \\
\hline
ResidentIncidents & Reporte y seguimiento de reclamos e incidencias. & CU\_17 \\
\hline
ResidentAmenities & Consulta de disponibilidad y reserva de áreas comunes. & CU\_05 \\
\hline
ResidentCartola & Estado de cuenta del residente (gastos comunes). & CU\_04 \\
\hline
ResidentChargesDetail & Detalle de cargos por período. & CU\_04 \\
\hline
ResidentPaymentFlow & Flujo de pago de gastos comunes. & CU\_04 \\
\hline
ResidentParcels & Visualización de encomiendas pendientes y retiradas. & CU\_02 \\
\hline
ResidentPublications & Foro comunitario con categorías e hilos. & CU\_16 \\
\hline
ResidentLibrary & Biblioteca de documentos comunitarios. & CU\_22 \\
\hline
ResidentMarketplace & Marketplace de artículos entre vecinos. & CU\_21 \\
\hline
ResidentMarketplaceCreate & Publicación de nuevo artículo en marketplace. & CU\_21 \\
\hline
ChatHub & Hub de conversaciones en tiempo real (WebSocket). & CU\_20 \\
\hline
ChatRoom & Sala de chat individual con historial. & CU\_20 \\
\hline
ResidentVotaciones & Listado de votaciones activas y cerradas. & CU\_10 \\
\hline
ResidentVoteDetail & Detalle de votación con opciones y resultados. & CU\_10 \\
\hline
ResidentTasks & Visualización de tareas asignadas (rol Personal). & CU\_13 \\
\hline
ResidentContacts & Gestión de contactos frecuentes para visitas. & CU\_01 \\

\end{longtable}

% --- Vistas Proveedor ---
\begin{longtable}{|p{0.28\textwidth}|p{0.42\textwidth}|p{0.22\textwidth}|}
\caption{Vistas del Proveedor}\label{tab:vistas-proveedor}\\
\hline
\textbf{Vista} & \textbf{Descripción} & \textbf{CUs relacionados} \\
\hline
\endfirsthead
\hline \textbf{Vista} & \textbf{Descripción} & \textbf{CUs relacionados} \\ \hline
\endhead
\hline \endfoot

ProviderPortal & Portal principal del proveedor con órdenes asignadas. & CU\_27 \\
\hline
ProviderServiceOrders & Detalle de órdenes con flujo de aceptación y cotizaciones. & CU\_27 \\

\end{longtable}

% --- Vistas Especiales ---
\begin{longtable}{|p{0.28\textwidth}|p{0.42\textwidth}|p{0.22\textwidth}|}
\caption{Vistas especiales y transversales}\label{tab:vistas-especiales}\\
\hline
\textbf{Vista} & \textbf{Descripción} & \textbf{CUs relacionados} \\
\hline
\endfirsthead
\hline \textbf{Vista} & \textbf{Descripción} & \textbf{CUs relacionados} \\ \hline
\endhead
\hline \endfoot

NotificationCenter & Centro de notificaciones con lectura y preferencias. & CU\_23 \\
\hline
CommunitySelector & Selector de comunidad activa (multi-comunidad). & CU\_24 \\
\hline
UserTypeSelector & Selección de tipo de usuario en registro de proveedores. & CU\_25 \\
\hline
NotFound (404) & Página de error para rutas inexistentes. & --- \\

\end{longtable}

% === PLACEHOLDER: Reemplazar con imagen generada ===
% PROMPT para generador de imágenes (Gemini/GPT/DALL-E):
%
% "Genera un MAPA DE NAVEGACIÓN (sitemap) para una aplicación web llamada DOMU.
%
% NOTACIÓN:
% - VISTAS/PANTALLAS: rectángulos pequeños con bordes redondeados, nombre de la vista dentro.
% - ROLES: 5 franjas horizontales (swimlanes) separadas por líneas, cada una con el nombre del rol a la izquierda.
% - CONEXIONES: líneas finas conectando vistas relacionadas (ej: Cartola → ChargesDetail → PaymentFlow).
% - VISTA COMPARTIDA: el rectángulo 'NotificationCenter' cruza todas las franjas con borde punteado.
%
% FRANJAS (de arriba a abajo):
% 1. 'Público' (fondo gris claro, 7 cajas): Home, Login, Register, ForgotPassword, ConfirmEmail, About, RegisterByInvite
% 2. 'Administrador' (fondo naranja pálido, 13 cajas): Dashboard, CreateUser, IncidentsBoard, IncidentStats,
%    Residents, HousingUnits, Amenities, CommonExpenses, Parcels, Tasks, Staff, Providers, ServiceOrders
% 3. 'Conserje' (fondo verde pálido, 3 cajas): VisitRegistration, VisitLog, Parcels
% 4. 'Residente' (fondo azul pálido, 18 cajas): Profile, Visits, Incidents, Amenities, Cartola,
%    ChargesDetail, PaymentFlow, Parcels, Publications, Library, Marketplace, MarketplaceCreate,
%    ChatHub, ChatRoom, Votaciones, VoteDetail, Tasks, Contacts
% 5. 'Proveedor' (fondo morado pálido, 2 cajas): ProviderPortal, ProviderServiceOrders
%
% Las cajas dentro de cada franja se organizan en filas. Conectar con flechas finas:
% Login → (cualquier dashboard según rol). Cartola → ChargesDetail → PaymentFlow.
% ChatHub → ChatRoom. Marketplace → MarketplaceCreate. Votaciones → VoteDetail.
%
% ESTILO: Fondo blanco general, franjas con colores pastel muy suaves. Bordes de cajas en gris oscuro.
% Sin sombras ni gradientes. Tipografía sans-serif pequeña. Profesional para tesis universitaria.
% Etiquetas en español. Formato: PNG 300dpi orientación horizontal (landscape)."
\begin{figure}[H]
  \centering
  \fbox{\parbox{0.85\textwidth}{\centering\vspace{3cm}
  \textit{[Insertar diagrama: Mapa de navegación del sistema por rol --- 54 vistas organizadas por actor]}
  \vspace{3cm}}}
  \caption{Mapa de navegación del sistema DOMU segmentado por rol de usuario.}
  \label{fig:mapa-navegacion-roles}
\end{figure}

\clearpage

\section{Diseño de arquitectura}
La solución adopta un patrón de arquitectura \textbf{Cliente-Servidor} con una API RESTful como medio principal de comunicación, complementada por canales WebSocket para eventos en tiempo real. Esta separación permite que el \textit{frontend} y el \textit{backend} evolucionen, se desarrollen y se desplieguen de manera independiente, favoreciendo la mantenibilidad y la escalabilidad horizontal.

La filosofía de diseño del \textit{backend} se fundamenta en tres principios: (i)~\textbf{separación de responsabilidades} mediante una arquitectura estricta de tres capas (\textit{Controller}~$\rightarrow$~\textit{Service}~$\rightarrow$~\textit{Repository}); (ii)~\textbf{inversión de dependencias} a través del contenedor de inyección Google~Guice, que desacopla la construcción de objetos de su uso; y (iii)~\textbf{aislamiento por comunidad} (\textit{multi-tenancy} lógico) mediante el \textit{header} HTTP \texttt{X-Building-Id}, que permite que una misma instancia del servidor atienda múltiples comunidades de forma segura.

La Tabla~\ref{tab:arquitectura-componentes} detalla los componentes principales y la tecnología que los sustenta.

\begin{longtable}{|p{0.28\textwidth}|p{0.64\textwidth}|}
\caption{Componentes arquitectónicos del sistema DOMU}\label{tab:arquitectura-componentes}\\
\hline
\textbf{Componente} & \textbf{Tecnología y Responsabilidad} \\
\hline
\endfirsthead
\hline
\textbf{Componente} & \textbf{Tecnología y Responsabilidad} \\
\hline
\endhead
\hline
\endfoot

Backend (Servidor) &
Java~21 (LTS) sobre JVM. Framework Javalin~6.7 para \textit{endpoints} HTTP y WebSocket. Arquitectura interna \textit{Controller-Service-Repository} con inyección de dependencias vía Google~Guice~7.0. \\
\hline

Frontend (Cliente) &
SPA desarrollada con React~19.1 y React~Router~7.9. Empaquetado con Vite~7.1, que optimiza tiempos de carga mediante \textit{code splitting} y carga perezosa (\textit{lazy loading}). \\
\hline

Base de Datos &
MySQL~8 gestionada con HikariCP~5.1 como \textit{pool} de conexiones. Esquema versionado mediante Flyway~10.15 (35~migraciones). \\
\hline

Autenticación &
Tokens JWT (java-jwt~4.4) con firma HMAC. Contraseñas almacenadas con \textit{hash} BCrypt. Middleware \texttt{AuthenticationHandler} valida cada petición protegida. \\
\hline

WebSocket &
Dos canales persistentes: \texttt{/ws/chat} (mensajería entre usuarios) y \texttt{/ws/notifications} (alertas en tiempo real). Autenticación por token en la conexión inicial. \\
\hline

Almacenamiento en la nube &
\texttt{BoxStorageService} como capa de abstracción sobre Box~API (SDK~5.2). Servicios especializados delegan la subida de avatares, evidencias y documentos. \\
\hline

Generación de documentos &
OpenPDF~1.3.39 para la generación de boletas de gasto común, comprobantes de pago y estados de cuenta en formato PDF. \\
\hline

Correo electrónico &
Jakarta~Mail~2.0.1 para notificaciones transaccionales: confirmación de registro, aprobación de comunidades e invitaciones de administrador. \\
\hline

Serialización &
Jackson~2.17 con módulo JSR-310 para el manejo nativo de \texttt{LocalDate} y \texttt{LocalDateTime}. Configurado globalmente mediante \texttt{AppConfig.CustomJsonMapper}. \\
\hline

Multi-tenancy &
\textit{Header} \texttt{X-Building-Id} enviado por el cliente en cada petición. El \textit{middleware} inyecta el identificador de comunidad, y los repositorios filtran todas las consultas por este valor. \\

\end{longtable}

Los patrones de diseño empleados incluyen: \textbf{Repository} para encapsular el acceso a datos, \textbf{DTO} (\textit{Data Transfer Object}) para la transferencia entre capas, \textbf{Middleware} para la validación transversal de autenticación y multi-tenancy, e \textbf{Inyección de Dependencias} para el desacoplamiento y la testabilidad del código.

\subsection{Configuración Inicial del Entorno}
Para garantizar la reproducibilidad del entorno de desarrollo y la correcta ejecución del sistema, se definieron configuraciones específicas tanto para el servidor como para el cliente, eliminando dependencias obsoletas (como \texttt{pom.xml} en favor de Gradle) y estandarizando las versiones de las herramientas.

\subsubsection*{Configuración del Backend}
El proyecto \texttt{domu-backend} se gestiona mediante \textbf{Gradle~8} (usando Kotlin~DSL en \texttt{build.gradle.kts}). Las versiones de todas las dependencias se centralizan en el catálogo \texttt{gradle/libs.versions.toml}, facilitando la actualización coordinada. El empaquetado final se realiza con el \textit{plugin} \textbf{Shadow~JAR}, generando un único archivo ejecutable (\textit{fat~JAR}) que incluye todas las dependencias transitivas.

La Tabla~\ref{tab:deps-backend} detalla las dependencias de producción del backend.

\begin{longtable}{|p{0.30\textwidth}|p{0.12\textwidth}|p{0.50\textwidth}|}
\caption{Dependencias de producción del backend}\label{tab:deps-backend}\\
\hline
\textbf{Librería} & \textbf{Versión} & \textbf{Propósito} \\
\hline
\endfirsthead
\hline
\textbf{Librería} & \textbf{Versión} & \textbf{Propósito} \\
\hline
\endhead
\hline
\endfoot

Javalin & 6.7.0 & Framework web ligero para HTTP y WebSocket. \\
\hline
MySQL Connector/J & 9.0.0 & Driver JDBC para MySQL~8. \\
\hline
HikariCP & 5.1.0 & Pool de conexiones JDBC de alto rendimiento. \\
\hline
Flyway Core + MySQL & 10.15.0 & Migraciones SQL versionadas al arranque. \\
\hline
Google Guice & 7.0.0 & Contenedor de inyección de dependencias. \\
\hline
java-jwt (Auth0) & 4.4.0 & Generación y validación de tokens JWT. \\
\hline
jBCrypt & 0.4 & \textit{Hashing} seguro de contraseñas. \\
\hline
Jackson JSR-310 & 2.17.2 & Serialización JSON con soporte \texttt{LocalDate}/\texttt{LocalDateTime}. \\
\hline
Box Java SDK & 5.2.0 & Almacenamiento de archivos en Box. \\
\hline
Google Cloud Storage & 2.36.1 & Almacenamiento alternativo en GCS. \\
\hline
OpenPDF & 1.3.39 & Generación de documentos PDF (boletas, estados de cuenta). \\
\hline
Jakarta Mail & 2.0.1 & Envío de correos transaccionales (SMTP). \\
\hline
Thumbnailator & 0.4.20 & Redimensionamiento de imágenes (avatares, evidencias). \\
\hline
Log4j~2 (Core + SLF4J) & 2.23.1 & Registro de eventos (\textit{logging}) estructurado. \\
\hline
SLF4J API & 2.0.13 & Fachada de \textit{logging}. \\
\hline
Javalin OpenAPI Plugin & 6.7.0 & Generación automática de especificación OpenAPI~3 desde anotaciones. \\
\hline
Javalin Swagger Plugin & 6.7.0 & Interfaz Swagger~UI integrada para exploración interactiva de la API. \\

\end{longtable}

La configuración en tiempo de ejecución se maneja a través de variables de entorno, priorizando la seguridad (no \textit{hardcodear} credenciales). Las variables críticas definidas en \texttt{AppConfig} son: \texttt{DB\_URI}, \texttt{DB\_USER}, \texttt{DB\_PASSWORD} para la conexión JDBC; \texttt{JWT\_SECRET} para la firma de tokens; \texttt{BOX\_DEVELOPER\_TOKEN} y \texttt{BOX\_ROOT\_FOLDER\_ID} para el almacenamiento; \texttt{MAIL\_HOST}, \texttt{MAIL\_USER} y \texttt{MAIL\_PASSWORD} para el correo; y \texttt{APP\_SERVER\_PORT} para el puerto de despliegue (por defecto \textbf{8080}).

\subsubsection*{Configuración del Frontend}
El proyecto \texttt{domu-frontend} utiliza \textbf{Node.js} (v18+) como entorno de ejecución y \textbf{npm} para la gestión de paquetes. La configuración base se encuentra en \texttt{vite.config.js} y \texttt{package.json}. En entorno de desarrollo, Vite configura un \textit{proxy} que redirige todas las peticiones a \texttt{/api} hacia \texttt{http://localhost:8080}, eliminando problemas de CORS sin necesidad de configuración adicional en el backend.

La Tabla~\ref{tab:deps-frontend} detalla las dependencias del frontend.

\begin{longtable}{|p{0.32\textwidth}|p{0.12\textwidth}|p{0.48\textwidth}|}
\caption{Dependencias del frontend}\label{tab:deps-frontend}\\
\hline
\textbf{Librería} & \textbf{Versión} & \textbf{Propósito} \\
\hline
\endfirsthead
\hline
\textbf{Librería} & \textbf{Versión} & \textbf{Propósito} \\
\hline
\endhead
\hline
\endfoot

\multicolumn{3}{|l|}{\textit{\textbf{Dependencias de producción}}} \\
\hline
React & 19.1.1 & Biblioteca de construcción de interfaces de usuario. \\
\hline
React DOM & 19.1.1 & Renderizado de componentes React en el DOM. \\
\hline
React Router DOM & 7.9.4 & Enrutamiento declarativo para la SPA. \\
\hline
TanStack React Query & 5.90.21 & Gestión de estado del servidor, caché y sincronización. \\
\hline
pigeon-maps & 0.21.6 & Componente de mapas interactivos (ubicación de comunidades). \\
\hline

\multicolumn{3}{|l|}{\textit{\textbf{Dependencias de desarrollo}}} \\
\hline
Vite & 7.1.7 & Empaquetador con HMR y \textit{code splitting}. \\
\hline
@vitejs/plugin-react & 5.0.4 & Integración de React con Vite (Fast Refresh). \\
\hline
vite-plugin-pwa & 1.1.0 & Soporte PWA con \textit{Service Worker} y manifiesto. \\
\hline
Sass & 1.97.3 & Preprocesador CSS con variables y \textit{mixins}. \\
\hline
ESLint & 9.36.0 & Análisis estático de código JavaScript/JSX. \\

\end{longtable}

Se utiliza un archivo \texttt{.env} (basado en \texttt{.env.example}) para definir \texttt{VITE\_API\_BASE\_URL}, permitiendo cambiar el apuntamiento de la API entre el entorno de desarrollo local (proxy) y producción sin modificar el código fuente.

\subsection{Arquitectura de Implementación Final}

Esta sección describe cómo interactúan los componentes arquitectónicos presentados anteriormente durante la ejecución real del sistema, abordando el flujo de una petición HTTP típica, los canales WebSocket, la estrategia de almacenamiento y el mecanismo de migraciones.

\subsubsection*{Flujo de una petición HTTP}

Cada solicitud del cliente recorre las siguientes etapas antes de devolver una respuesta:

\begin{enumerate}
    \item \textbf{Cliente (React)}: El componente React invoca una función del módulo de servicios API, que utiliza \texttt{fetch} con las cabeceras \texttt{Authorization: Bearer~<token>} y \texttt{X-Building-Id:~<id>}.
    \item \textbf{Proxy de desarrollo (Vite)}: En entorno local, Vite intercepta las rutas \texttt{/api/*} y las redirige a \texttt{http://localhost:8080}, eliminando restricciones de CORS.
    \item \textbf{Javalin (Router)}: El servidor recibe la petición y la enruta al \textit{handler} correspondiente según el método HTTP y la ruta registrada.
    \item \textbf{AuthenticationHandler (Middleware)}: Valida el token JWT, extrae el identificador de usuario y lo inyecta en el contexto de la petición. Las rutas públicas (login, registro, health) se excluyen de esta validación.
    \item \textbf{Controller}: Deserializa el cuerpo de la petición mediante Jackson, valida los parámetros de entrada y delega la lógica al \textit{Service} correspondiente.
    \item \textbf{Service}: Implementa la lógica de negocio, coordina múltiples repositorios cuando es necesario y aplica el filtrado por \texttt{buildingId} para garantizar el aislamiento multi-comunidad.
    \item \textbf{Repository}: Ejecuta las consultas SQL parametrizadas contra MySQL a través del \textit{pool} de conexiones HikariCP y devuelve las entidades al \textit{Service}.
    \item \textbf{Respuesta}: El \textit{Controller} serializa el resultado como JSON (con soporte JSR-310 para fechas) y lo devuelve al cliente con el código HTTP apropiado.
\end{enumerate}

\subsubsection*{Canales WebSocket}

El sistema mantiene dos canales WebSocket persistentes, ambos autenticados mediante token JWT enviado como parámetro en la URL de conexión:

\begin{itemize}
    \item \texttt{/ws/chat}: Gestionado por \texttt{ChatWebSocketHandler}. Permite el envío y recepción de mensajes en tiempo real entre usuarios de una misma comunidad. El \textit{handler} mantiene un mapa de sesiones activas por usuario para la entrega inmediata.
    \item \texttt{/ws/notifications}: Gestionado por \texttt{NotificationWebSocketHandler}. Emite alertas en tiempo real (nuevas encomiendas, incidentes asignados, pagos registrados, etc.) según las preferencias configuradas por cada usuario.
\end{itemize}

\subsubsection*{Estrategia de almacenamiento}

Los archivos binarios (avatares, evidencias fotográficas, documentos de biblioteca) se almacenan externamente a través de \texttt{BoxStorageService}, que encapsula la interacción con la API de Box. Los servicios especializados (\texttt{UserProfileService}, \texttt{ParcelService}, \texttt{LibraryService}) delegan las operaciones de carga y descarga a esta capa centralizada, manteniendo en la base de datos únicamente la referencia (\texttt{fileId}) al recurso remoto. Las imágenes se redimensionan automáticamente mediante Thumbnailator antes de su almacenamiento.

\subsubsection*{Migraciones de base de datos}

Al arranque del servidor, Flyway escanea el directorio \texttt{db/migration} y ejecuta secuencialmente las migraciones pendientes. Cada archivo sigue la convención \texttt{V\{NNN\}\_\_descripcion.sql}, donde el prefijo numérico garantiza el orden de aplicación. El historial de migraciones se registra en la tabla interna \texttt{flyway\_schema\_history}, lo que permite detectar modificaciones manuales no autorizadas y garantizar la consistencia entre entornos.

\subsubsection*{Documentación de la API (Swagger UI)}

La API REST se documenta de forma automática mediante el plugin comunitario \textbf{javalin-openapi}, que genera una especificación OpenAPI~3 a partir de anotaciones en el código fuente. La interfaz interactiva \textbf{Swagger~UI} se encuentra disponible en la ruta \texttt{/swagger}, permitiendo a los desarrolladores explorar los endpoints, probar peticiones y consultar los esquemas de entrada/salida sin herramientas externas. La especificación JSON se expone en \texttt{/openapi} para su consumo programático.

\subsubsection*{Serialización y manejo de fechas}

La serialización JSON se configura globalmente mediante \texttt{AppConfig.CustomJsonMapper}, que registra el módulo \texttt{JavaTimeModule} de Jackson y desactiva la escritura de fechas como \textit{timestamps} numéricos. Esto permite que los tipos \texttt{LocalDate} y \texttt{LocalDateTime} se representen como cadenas ISO-8601 (e.g., \texttt{"2025-03-15"}, \texttt{"2025-03-15T14:30:00"}), facilitando su procesamiento tanto en el frontend como en integraciones externas.
